\documentclass[11pt,a4paper]{article}
\usepackage[utf8]{inputenc}
\usepackage[T1]{fontenc}
\usepackage{amsmath,amsfonts,amssymb}
\usepackage{graphicx}
\usepackage{hyperref}
\usepackage{listings}
\usepackage{color}
\usepackage{booktabs}
\usepackage{float}
\usepackage{geometry}
\geometry{margin=1in}
\definecolor{zenblue}{RGB}{41,121,255}
\definecolor{zengreen}{RGB}{52,199,89}
\definecolor{codegray}{RGB}{240,240,240}
\hypersetup{colorlinks=true,linkcolor=zenblue,urlcolor=zenblue,citecolor=zenblue}

\lstset{
  backgroundcolor=\color{codegray},
  basicstyle=\ttfamily\small,
  breaklines=true,
  frame=single,
  language=Python
}

\title{\textbf{Zen-Code: Code Completion and Intelligence at 14B Scale}\\
\large Technical Report v2024.12}
\author{Zen LM Research Team\\
\texttt{research@zenlm.org}}
\date{December 2024}

\begin{document}
\maketitle

\begin{abstract}
We present \textbf{Zen-Code}, a 14 billion parameter language model optimized for code
generation, completion, and understanding across 80$+$ programming languages. Built on the
Zen MoDE (Mixture of Distilled Experts) architecture, Zen-Code is pretrained on 2T tokens of
high-quality code with fill-in-the-middle (FIM) objectives and fine-tuned on repository-scale
context tasks. Zen-Code achieves HumanEval 87.2\%, MBPP 82.3\%, SWE-bench 28.4\%, and
RepoBench 0.834, establishing it as the capable code-specialist in the Zen family. The model
supports long repository-level context up to 64K tokens, multi-file awareness, and standard
IDE integration formats including LSP-compatible completions.
\end{abstract}

\tableofcontents
\newpage

%% ─────────────────────────────────────────────────────────────────────────────
\section{Introduction}

Code generation has emerged as one of the most commercially significant applications of large
language models. From IDE autocomplete to automated pull request generation, models that reliably
understand and generate code at the repository scale provide substantial productivity gains
\cite{chen2021codex, nijkamp2023codegen2}. General-purpose language models, while capable of
code generation, are not optimized for the specific properties of software: multi-file context,
precise syntax requirements, test-driven correctness, and sub-token completion granularity.

Zen-Code addresses this gap with a model purpose-built for code. Key contributions:

\begin{itemize}
  \item A 14B parameter model pretrained on 2T code tokens from 80$+$ languages, with an
        additional 500B tokens of technical documentation and Stack Exchange content.
  \item Fill-in-the-middle (FIM) training enabling prefix-suffix-middle completion, the
        format required by production IDE integrations.
  \item Repository-level context training on 64K token windows with structured file-path
        and symbol-table prefixes.
  \item Post-training on instruction-following code tasks including debugging, refactoring,
        test generation, and code review.
  \item Benchmark results: HumanEval 87.2\%, MBPP 82.3\%, SWE-bench 28.4\%, RepoBench 0.834.
\end{itemize}

%% ─────────────────────────────────────────────────────────────────────────────
\section{Architecture}

\subsection{Model Configuration}

Zen-Code uses the Zen MoDE architecture scaled to 14B parameters, a tier chosen to balance
code reasoning capability with inference cost on developer hardware (single A100 or 2$\times$
RTX 4090 in BF16).

\begin{table}[H]
\centering
\caption{Zen-Code architecture hyperparameters.}
\label{tab:arch}
\begin{tabular}{lc}
\toprule
\textbf{Hyperparameter} & \textbf{Value} \\
\midrule
Parameters (total)           & 14.4B \\
Layers                       & 40 \\
Attention heads              & 40 \\
KV heads (GQA)               & 8 \\
Hidden dimension             & 5120 \\
FFN intermediate dimension   & 13{,}696 \\
Vocabulary size              & 151{,}936 \\
Context length (training)    & 65{,}536 \\
Position encoding            & RoPE ($\theta = 2{,}000{,}000$) \\
Activation                   & SiLU \\
Normalization                & RMSNorm \\
FIM tokens                   & \texttt{<|fim\_prefix|>}, \texttt{<|fim\_middle|>}, \texttt{<|fim\_suffix|>} \\
\bottomrule
\end{tabular}
\end{table}

\subsection{Fill-in-the-Middle Token Scheme}

Zen-Code supports fill-in-the-middle (FIM) completion \cite{bavarian2022fim}, the standard
format for IDE inline completion. The three FIM sentinel tokens partition the input:

\begin{lstlisting}[language={},caption={FIM token format.}]
<|fim_prefix|>[prefix code]<|fim_suffix|>[suffix code]<|fim_middle|>
\end{lstlisting}

During training, 50\% of code samples are converted to FIM format using the PSM
(prefix-suffix-middle) or SPM (suffix-prefix-middle) transformation drawn with equal
probability, following the recommendation of \cite{bavarian2022fim}.

For the PSM transformation of a document $d$ split at position $s$:

\begin{equation}
  d_{\text{FIM}} = [\texttt{FP}] \cdot d_{[:s]} \cdot [\texttt{FS}] \cdot d_{[s:]} \cdot [\texttt{FM}] \cdot d_{\text{middle}}
\end{equation}

where $\texttt{FP}, \texttt{FS}, \texttt{FM}$ are the prefix, suffix, and middle sentinels
respectively, and $d_{\text{middle}}$ is the held-out span the model must predict.

\subsection{Repository-Level Context Structure}

For long-context code tasks, Zen-Code accepts structured repository context using a
lightweight file-header format:

\begin{lstlisting}[language={},caption={Repository context format.}]
# File: src/utils/parser.py
[file content]

# File: src/models/base.py
[file content]

# File: [target file, completion requested here]
[prefix]<|fim_suffix|>[suffix]<|fim_middle|>
\end{lstlisting}

This format is injected by IDE extensions and agents to provide cross-file context.
The 64K context window accommodates typical repository relevant-file sets (10--50 files
of 500--1000 lines each).

%% ─────────────────────────────────────────────────────────────────────────────
\section{Training Methodology}

\subsection{Pretraining Data}

Zen-Code's 2.5T-token pretraining corpus is heavily code-dominated.

\begin{table}[H]
\centering
\caption{Zen-Code pretraining data composition.}
\label{tab:data}
\begin{tabular}{lcc}
\toprule
\textbf{Domain} & \textbf{Tokens (B)} & \textbf{Fraction} \\
\midrule
Source code (80$+$ languages)     & 1{,}500 & 60.0\% \\
Technical documentation           &   375 & 15.0\% \\
Stack Overflow / Q\&A             &   250 & 10.0\% \\
Research papers (CS/ML)           &   125 &  5.0\% \\
General web (filtered)            &   125 &  5.0\% \\
Formal specs (Lean, Coq, TLA$+$) &    75 &  3.0\% \\
Synthetic code (unit tests)       &    50 &  2.0\% \\
\midrule
Total                             & 2{,}500 & 100.0\% \\
\bottomrule
\end{tabular}
\end{table}

\subsection{Per-Language Distribution}

\begin{table}[H]
\centering
\caption{Top programming languages by token count in pretraining corpus.}
\label{tab:langs}
\begin{tabular}{lcc}
\toprule
\textbf{Language} & \textbf{Tokens (B)} & \textbf{Share of code data} \\
\midrule
Python     & 360 & 24.0\% \\
JavaScript & 240 & 16.0\% \\
TypeScript & 150 &  10.0\% \\
Java       & 135 &  9.0\% \\
C/C$+$$+$  & 120 &  8.0\% \\
Rust       &  90 &  6.0\% \\
Go         &  75 &  5.0\% \\
Ruby       &  45 &  3.0\% \\
PHP        &  45 &  3.0\% \\
Shell/Bash &  45 &  3.0\% \\
Other (70$+$) & 195 & 13.0\% \\
\midrule
Total      & 1{,}500 & 100.0\% \\
\bottomrule
\end{tabular}
\end{table}

Data deduplication uses a two-stage process: exact SHA-256 deduplication at the file level,
followed by MinHash near-deduplication with 10-shingling at LSH threshold 0.8 to remove
forked repositories and boilerplate.

\subsection{Pretraining Procedure}

Zen-Code is trained from scratch (not initialized from Zen base) to allow the training
distribution to be fully code-optimized without general text forgetting constraints.

\begin{itemize}
  \item Optimizer: AdamW, LR $3 \times 10^{-4}$ $\to$ $3 \times 10^{-5}$ (cosine)
  \item Warm-up: 1000 steps
  \item Weight decay: 0.1
  \item Batch size: 4M tokens
  \item FIM rate: 50\% of code samples
  \item Precision: BF16
  \item Training duration: $\approx$625K steps, $\approx$512 H100 GPUs
\end{itemize}

\subsection{Repository-Level Fine-Tuning}

After base pretraining, Zen-Code undergoes a 50B-token repository-level fine-tuning phase
using full repository snapshots. Repositories are sampled from GitHub with permissive licenses,
filtered for non-trivial size (>10 files, >1K lines). Files within each repository are packed
into 64K-token windows following a topological ordering based on import dependencies.

\subsection{Instruction Fine-Tuning}

A final SFT phase on 800K instruction-code pairs covering:

\begin{itemize}
  \item Code generation from natural language descriptions
  \item Bug finding and fixing (with error message context)
  \item Code explanation and documentation generation
  \item Unit test generation from function signatures
  \item Code refactoring and style improvement
  \item API usage from documentation
\end{itemize}

SFT learning rate: $1 \times 10^{-5}$, 2 epochs.

%% ─────────────────────────────────────────────────────────────────────────────
\section{Evaluation}

\subsection{Code Generation Benchmarks}

\begin{table}[H]
\centering
\caption{Code generation benchmark results.}
\label{tab:benchmarks}
\begin{tabular}{lcccc}
\toprule
\textbf{Benchmark} & \textbf{Zen-Code (14B)} & \textbf{Comp.\ A (15B)} & \textbf{Comp.\ B (13B)} & \textbf{Comp.\ C (7B)} \\
\midrule
HumanEval (pass@1)      & \textbf{87.2} & 85.1 & 83.6 & 76.3 \\
HumanEval$+$ (pass@1)   & \textbf{82.4} & 80.3 & 78.7 & 71.2 \\
MBPP (pass@1)           & \textbf{82.3} & 80.1 & 78.4 & 70.8 \\
MBPP$+$ (pass@1)        & \textbf{76.8} & 74.3 & 72.1 & 64.3 \\
LiveCodeBench (3-month) & \textbf{54.2} & 51.7 & 49.3 & 41.8 \\
\midrule
SWE-bench Verified      & \textbf{28.4} & 25.3 & 22.8 & 14.6 \\
RepoBench (retrieval)   & \textbf{0.834} & 0.812 & 0.798 & 0.741 \\
\midrule
MultiPL-E (avg 12 langs)& \textbf{72.4} & 70.8 & 68.3 & 61.2 \\
\bottomrule
\end{tabular}
\end{table}

\subsection{Multilingual Code Evaluation}

\begin{table}[H]
\centering
\caption{MultiPL-E pass@1 by programming language.}
\label{tab:multilingual_code}
\begin{tabular}{lcc}
\toprule
\textbf{Language} & \textbf{Zen-Code (14B)} & \textbf{Competitor (15B)} \\
\midrule
Python      & 87.2 & 85.1 \\
JavaScript  & 78.6 & 76.4 \\
TypeScript  & 76.4 & 74.1 \\
Java        & 74.3 & 72.8 \\
C$+$$+$     & 70.8 & 68.3 \\
Rust        & 68.4 & 65.7 \\
Go          & 73.1 & 71.4 \\
Ruby        & 65.3 & 62.8 \\
PHP         & 63.7 & 61.4 \\
Shell       & 58.4 & 55.9 \\
\bottomrule
\end{tabular}
\end{table}

\subsection{Fill-in-the-Middle Evaluation}

\begin{table}[H]
\centering
\caption{FIM completion accuracy on HumanEval-Infilling.}
\label{tab:fim}
\begin{tabular}{lcc}
\toprule
\textbf{Metric} & \textbf{Zen-Code (14B)} & \textbf{Competitor (15B)} \\
\midrule
Single-line (exact match) & 82.3 & 79.1 \\
Multi-line (exact match)  & 61.4 & 57.8 \\
Single-line (functional)  & 91.7 & 89.4 \\
Multi-line (functional)   & 74.6 & 71.3 \\
\bottomrule
\end{tabular}
\end{table}

\subsection{Repository-Level Evaluation}

SWE-bench Verified requires understanding a full repository, identifying which files to
modify, and generating a correct patch that passes the repository's test suite.

\begin{table}[H]
\centering
\caption{SWE-bench Verified resolution rates by issue category.}
\label{tab:swebench}
\begin{tabular}{lc}
\toprule
\textbf{Category} & \textbf{Zen-Code (14B)} \\
\midrule
Logic errors        & 34.2\% \\
Type errors         & 41.8\% \\
Off-by-one          & 38.7\% \\
Missing handling    & 31.4\% \\
API usage errors    & 29.6\% \\
Regression fixes    & 18.3\% \\
\midrule
Overall (verified)  & 28.4\% \\
\bottomrule
\end{tabular}
\end{table}

%% ─────────────────────────────────────────────────────────────────────────────
\section{IDE Integration}

Zen-Code is designed for IDE integration via a thin inference server exposing an
OpenAI-compatible completions API with FIM support:

\begin{lstlisting}[language=Python, caption={Example FIM completion request.}]
import openai

client = openai.OpenAI(
    base_url="https://api.zenlm.org/v1",
    api_key="<zen-api-key>"
)

response = client.completions.create(
    model="zen-code-14b",
    prompt="<|fim_prefix|>def merge_sorted_arrays(a, b):\n    "
           "<|fim_suffix|>\n    return result\n<|fim_middle|>",
    max_tokens=128,
    temperature=0.1,
    stop=["<|endoftext|>"]
)

print(response.choices[0].text)
\end{lstlisting}

Available IDE integrations include VS Code (via Continue and Tabby plugins), JetBrains
IDEs, Neovim (via coc.nvim), and Emacs (via copilot.el-compatible backends).

%% ─────────────────────────────────────────────────────────────────────────────
\section{Ablation Studies}

\subsection{Effect of FIM Training Rate}

\begin{table}[H]
\centering
\caption{HumanEval-Infilling (single-line) vs.\ FIM training fraction.}
\label{tab:fim_ablation}
\begin{tabular}{lcc}
\toprule
\textbf{FIM rate} & \textbf{HumanEval pass@1} & \textbf{FIM exact match (single)} \\
\midrule
0\%  (no FIM)  & 86.8 & 41.2 \\
25\%           & 85.9 & 74.3 \\
50\%           & 87.2 & 82.3 \\
75\%           & 84.3 & 83.1 \\
\bottomrule
\end{tabular}
\end{table}

50\% FIM rate provides the optimal balance: full left-to-right completion performance is
maintained (87.2\% vs.\ 86.8\% baseline) while infilling performance improves dramatically.
At 75\% FIM the model under-trains on left-to-right generation, reducing HumanEval by 2.9pp.

\subsection{Effect of Repository Fine-Tuning}

\begin{table}[H]
\centering
\caption{Impact of repository-level fine-tuning on long-context tasks.}
\label{tab:repo_ablation}
\begin{tabular}{lcc}
\toprule
\textbf{Model} & \textbf{RepoBench} & \textbf{SWE-bench Verified} \\
\midrule
Base (no repo FT)         & 0.761 & 19.3 \\
$+$ Repo FT (8K context)  & 0.803 & 23.7 \\
$+$ Repo FT (64K context) & 0.834 & 28.4 \\
\bottomrule
\end{tabular}
\end{table}

%% ─────────────────────────────────────────────────────────────────────────────
\section{Related Work}

Code-specialized language models trace to Codex \cite{chen2021codex}, which demonstrated
that training on large code corpora enables competitive function-level generation. CodeGen
\cite{nijkamp2023codegen2} and its successors explored scaling laws for code. StarCoder
\cite{li2023starcoder} trained on The Stack, a permissively licensed code corpus with strong
deduplication. Fill-in-the-middle training was introduced in \cite{bavarian2022fim} and
adopted by most subsequent code models. SWE-bench \cite{jimenez2023swe} introduced
repository-level task solving as a benchmark; the verified subset reduces noise in evaluation.
RepoBench \cite{liu2023repobench} tests long-context retrieval within repository structure.

%% ─────────────────────────────────────────────────────────────────────────────
\section{Limitations}

Zen-Code's SWE-bench performance of 28.4\% indicates that most real-world repository-level
issues remain unsolved in single-pass generation without agentic scaffolding. The model does
not have access to runtime execution during generation; code correctness is assessed only at
test time. Languages with fewer than 10B training tokens show substantially weaker performance
than top-tier languages. Security-sensitive code generation (cryptography, access control)
should be reviewed by experts regardless of benchmark scores.

%% ─────────────────────────────────────────────────────────────────────────────
\section{Conclusion}

Zen-Code provides a purpose-built code intelligence model at 14B parameters, delivering
HumanEval 87.2\%, MBPP 82.3\%, SWE-bench 28.4\%, and RepoBench 0.834 through a combination
of code-dominated pretraining, fill-in-the-middle objectives, and repository-scale context
fine-tuning. The model's 64K context window, FIM support, and IDE-compatible API make it
a practical choice for developer tooling deployment. Zen-Code is the foundational code model
in the Zen family, with Zen-Coder-Flash providing a distilled fast-inference variant for
latency-critical IDE autocomplete applications.

%% ─────────────────────────────────────────────────────────────────────────────
\begin{thebibliography}{99}

\bibitem{chen2021codex}
M.~Chen et al., ``Evaluating Large Language Models Trained on Code,''
\textit{arXiv:2107.03374}, 2021.

\bibitem{nijkamp2023codegen2}
E.~Nijkamp et al., ``CodeGen2: Lessons for Training LLMs on Programming and Natural Languages,''
\textit{ICLR}, 2023.

\bibitem{bavarian2022fim}
M.~Bavarian et al., ``Efficient Training of Language Models to Fill in the Middle,''
\textit{arXiv:2207.14255}, 2022.

\bibitem{li2023starcoder}
R.~Li et al., ``StarCoder: May the Source Be with You!'' \textit{arXiv:2305.06161}, 2023.

\bibitem{jimenez2023swe}
C.~Jimenez et al., ``SWE-bench: Can Language Models Resolve Real-World GitHub Issues?''
\textit{arXiv:2310.06770}, 2023.

\bibitem{liu2023repobench}
T.~Liu et al., ``RepoBench: Benchmarking Repository-Level Code Auto-Completion Systems,''
\textit{arXiv:2306.03091}, 2023.

\bibitem{loshchilov2019decoupled}
I.~Loshchilov and F.~Hutter, ``Decoupled Weight Decay Regularization,'' \textit{ICLR}, 2019.

\bibitem{broder1997minwise}
A.~Broder, ``On the resemblance and containment of documents,'' \textit{Sequences}, 1997.

\end{thebibliography}

\end{document}
